\section{Logical base and the meaning of an embedded proof}

\subsection{The repository's ordinary presentation of iPRA}

The checked backend uses the following objects.

\begin{itemize}
\item \(\PR\,k\) is the inductively generated type of primitive-recursive
      function symbols of arity \(k\), with zero, successor, projections,
      composition, and primitive recursion.
\item \(\PTm\,n\) is the first-order term language with \(n\) free variables,
      formed by variables and applications of \(\PR\) symbols.
\item \(\PRAAx\) contains the defining equations for every primitive-recursive
      symbol, successor separation and injectivity, and induction for
      quantifier-free formulas.
\item \(\Prf\,\PRAAx\,\Gamma\,\varphi\), written
      \(\Gamma\vdash[\PRAAx]\varphi\), is intuitionistic natural deduction
      with equality, substitution, and open axioms.
\end{itemize}

The crucial repository theorem is
\[
(\Gamma\vdash[\PRAAx]\varphi)
\quad\Longleftrightarrow\quad
\NPrf(\Gamma^{\mathrm{nat}},\varphi^{\mathrm{nat}}),
\]
implemented as \code{pra\_iff\_native}.  It is a faithful transport theorem,
not a reflection principle from arbitrary Lean propositions.

\subsection{Schemas versus single first-order theorems}

A fixed primitive-recursive evaluator \(a:\N^{r+1}\to\N\) can be represented
by one fixed \(\PR\) symbol.  A theorem may then be uniform in its ordinary
numerical parameter \(p\) and precision \(k\):
\[
  \forall p\,\forall k\;\Phi_a(p,k).
\]
This is a single first-order theorem for that fixed symbol.  A Lean theorem
polymorphic in a meta-level value \code{a : PR (r+1)} is instead a schema
generator: for each fixed symbol it produces a first-order formula and proof.
The distinction is benign and should be explicit.

What is not legitimate is to quantify over arbitrary program codes \(e\) and
pretend that there is a total primitive-recursive function
\(U(e,x)\) evaluating every primitive-recursive program.  The diagonal
argument in \cref{sec:barriers} rules this out.  A theorem about an external
algorithm must therefore use one of two interfaces:

\begin{enumerate}
\item its evaluator is fixed in the theorem's language; or
\item a finite computation trace is supplied and checked by a fixed verifier.
\end{enumerate}

\subsection{Definitional extensions are harmless when fixed}

The usual mathematical presentation of iPRA has a function symbol for every
primitive-recursive function.  The repository instead represents these
symbols by syntax trees in \(\PR\), with defining equations in \(\PRAAx\).
This means that introducing a named helper such as pairing, list lookup,
integer addition, rational comparison, bounded maximum, or a particular trace
runner is not an additional axiom.  One defines a fixed \(\PR\) term and uses
its generated defining equations.

For proof engineering it is useful to distinguish three levels.

\begin{description}[leftmargin=3.2cm,style=nextline]
\item[Meta-level definition.] A Lean function constructs a \(\PR\) syntax
      tree.  This is not yet a theorem.
\item[Object term.] Applying the resulting symbol gives a \(\PTm\) term.
\item[Object proof.] Theorems such as \code{PRA.comp}, \code{PRA.precZero},
      and \code{PRA.precSucc} expose the defining equations as
      \(\Prf\,\PRAAx\) derivations.
\end{description}

A formalization should avoid duplicating encodings.  The same Boolean, pair,
list, signed-integer, and positive-denominator rational codes should be reused
across all packets, because every additional representation creates a new
extensionality and conversion burden.

\subsection{Quantifier-free induction with quantified context}

A frequent misunderstanding is that quantifier-free induction forbids the use
of represented-object hypotheses containing quantifiers.  It does not.  The
restriction applies to the induction formula, not to the surrounding context.
For example, suppose the context contains a regularity hypothesis
\[
  H(p):\equiv \forall k\,\forall \ell\;
    \dist(a(p,k),a(p,\ell))\le 2^{-k}+2^{-\ell}.
\]
Inside the successor step one may instantiate \(H(p)\) at primitive-recursive
terms computed from the current finite state.  The induction predicate itself
can remain the equation
\[
  \mathsf{accumulator}(i,p)=1.
\]
Thus a global quantified contract may be consumed stage by stage while the
formula on which induction acts is quantifier free.

This observation is the logical basis of the finite query-table discipline.
A computation declares or records every sample it reads.  The global contract
is instantiated only at those samples, and the resulting finite facts are
folded into a Boolean table.  The induction invariant says that the prefix of
the table already constructed is valid.

\subsection{The proof-producing interface}

The correct reusable abstraction is not an ambient proposition asserting that
an algorithm is sound.  It is a structure whose soundness field is already an
object-language derivation.

At a schematic level, a proof-producing transducer from a contract \(A\) to a
contract \(B\) consists of:
\[
  \text{an output term }F(\bar x),
  \qquad
  A(\bar x)\ \vdash[\PRAAx]\ B(\bar x,F(\bar x)).
\]
Composition is object-language cut.  If
\[
 A\vdash B(F) \quad\text{and}\quad B(y)\vdash C(G(y)),
\]
then substitution followed by \code{Prf.trans} gives
\[
 A\vdash C(G(F)).
\]
This simple formulation prevents a whole class of category mistakes.  The
semantic development may be large, but its reusable interface ends in a term
of the repository's derivation type.

\subsection{What the generic framework may assume}

The core theorems below assume a small arithmetic library whose existence is
standard and whose completion is a finite formalization task:

\begin{enumerate}
\item primitive-recursive pairing and finite-list coding, with projection and
      length equations;
\item Boolean values \(0,1\), Boolean conjunction/disjunction/negation, and
      normalization lemmas;
\item zero test, equality test, order, truncated subtraction, and bounded
      search;
\item signed integers and rationals represented by normalized or guarded
      codes; and
\item dyadic radii and elementary error-budget arithmetic.
\end{enumerate}

Nothing in the master trace theorem depends on the details of these codes.
For a concrete Lean implementation, however, selecting one canonical library
before building domain-specific modules is the first milestone.
